\documentclass[11pt]{article}
\usepackage{mathblatt}
% gesetzt von werkzeuge/setzer.py; Lösungen nach bau/bauregeln.md „Lösungen“.
\newcommand{\kennung}{KUV}
\begin{document}
\pfheftstil
\fancyfoot[C]{{\footnotesize\color{mbgrau}\kennung\hfill\thepage}}
\pfheftkopf{Wahrscheinlichkeit einstufig}{Klasse 7 \textperiodcentered{} Lösungen}

\begin{pfloesung}{}
\lz{1b)}{$P = \frac{1}{6}$}{möglich 6 (1 bis 6), günstig 1: $P = \frac{1}{6}$}{}
\end{pfloesung}
\begin{pfloesung}{}
\lz{2.}{$P = \frac{7}{12}$}{möglich $4 + 7 + 1 = 12$ Kugeln, günstig 7: $P = \frac{7}{12}$}{}
\end{pfloesung}
\begin{pfloesung}{}
\lz{3.}{$P = \frac{1}{2}$}{möglich 6 Felder, günstig 3 Felder mit A: $P = \frac{3}{6} = \frac{1}{2}$. Es zählen die Felder, nicht die Buchstaben.}{}
\end{pfloesung}
\begin{pfloesung}{}
\lz{4.}{$P = \frac{1}{5}$}{möglich 20, günstig 5, 10, 15, 20 (4 Zettel): $P = \frac{4}{20} = \frac{1}{5}$}{}
\end{pfloesung}
\begin{pfloesung}{}
\lz{5.}{Aus dem ersten: $\frac{3}{8} > \frac{1}{3}$}{Aus dem ersten Beutel. Erster Beutel: $P = \frac{3}{8} = 0{,}375$; zweiter Beutel: $P = \frac{4}{12} = \frac{1}{3} \approx 0{,}333$. Im zweiten sind zwar mehr rote Kugeln, aber ihr Anteil ist kleiner.}{}
\end{pfloesung}
\begin{pfloesung}{}
\lz{6b)}{$0{,}6 = 60\,\%$}{$3 : 5 = 0{,}6 = 60\,\%$}{}
\end{pfloesung}
\begin{pfloesung}{}
\lz{7.}{$\frac{1}{10} = 10\,\%$}{alle $54 + 6 = 60$, günstig 6: $P = \frac{6}{60} = \frac{1}{10} = 10\,\%$ (nicht $\frac{6}{54}$)}{}
\end{pfloesung}
\begin{pfloesung}{}
\lz{8.}{$44\,\%$}{alle $11 + 6 + 8 = 25$, günstig 11: $P = \frac{11}{25} = 0{,}44 = 44\,\%$}{}
\end{pfloesung}
\begin{pfloesung}{}
\lz{9.}{$\frac{1}{8} = 12{,}5\,\%$}{alle $600 - 201 + 1 = 400$ Lose, günstig $250 - 201 + 1 = 50$: $P = \frac{50}{400} = \frac{1}{8} = 0{,}125 = 12{,}5\,\%$}{}
\end{pfloesung}
\begin{pfloesung}{}
\lz{10.}{Bude A: $25\,\%$ gegen $22\,\%$}{Bude A. A: alle $150 + 50 = 200$, $P = \frac{50}{200} = 25\,\%$; B: $P = \frac{66}{300} = 22\,\%$. Falsch wäre $\frac{50}{150}$.}{}
\end{pfloesung}
\begin{pfloesung}{}
\lz{11b)}{$0{,}7$ (also $70\,\%$)}{$1 - 0{,}3 = 0{,}7$}{}
\end{pfloesung}
\begin{pfloesung}{}
\lz{12.}{$P = \frac{5}{6}$}{$P(6) = \frac{1}{6}$; $P(\text{keine Sechs}) = 1 - \frac{1}{6} = \frac{5}{6}$}{}
\end{pfloesung}
\begin{pfloesung}{}
\lz{13.}{$P = \frac{2}{3}$}{alle 12; $P(\text{grün}) = \frac{4}{12}$; $P(\text{nicht grün}) = 1 - \frac{4}{12} = \frac{8}{12} = \frac{2}{3}$}{}
\end{pfloesung}
\begin{pfloesung}{}
\lz{14.}{$P = \frac{1}{2}$}{durch 3 oder durch 4 teilbar: 3, 4, 6, 8, 9, 12 – das sind 6 (die 12 nur einmal zählen); $P = 1 - \frac{6}{12} = \frac{6}{12} = \frac{1}{2}$}{}
\end{pfloesung}
\begin{pfloesung}{}
\lz{15.}{Nein: Ole $\frac{2}{3}$, Mia nur $\frac{1}{3}$}{Nein. $P(\text{Mia}) = P(1 \text{ oder } 2) = \frac{2}{6} = \frac{1}{3}$; $P(\text{Ole}) = 1 - \frac{1}{3} = \frac{2}{3}$. Ole gewinnt doppelt so oft wie Mia.}{}
\end{pfloesung}
\begin{pfloesung}{}
\lz{16b)}{$P = \frac{1}{3}$}{noch $10 - 1 = 9$ Kugeln, rot 3: $P = \frac{3}{9} = \frac{1}{3}$}{}
\end{pfloesung}
\begin{pfloesung}{}
\lz{17.}{$P = \frac{4}{11}$}{noch $12 - 1 = 11$ Kugeln, davon $5 - 1 = 4$ gelb: $P = \frac{4}{11}$}{}
\end{pfloesung}
\begin{pfloesung}{}
\lz{18.}{$P = \frac{3}{25} = 12\,\%$}{möglich sind nur noch die Nummern mit Endziffer 3: 103, 113, \ldots, 343 – das sind 25. Davon enden auf 13: 113, 213, 313, also 3. $P = \frac{3}{25} = 0{,}12 = 12\,\%$ (nicht $\frac{3}{250}$)}{}
\end{pfloesung}
\begin{pfloesung}{}
\lz{19.}{Nein: $\frac{4}{50} = 8\,\%$}{Nein. Lose mit Endziffer 7: 307, 317, \ldots, 797 – das sind 50. Davon enden auf 27: 327, 427, 527, 627, 727; ohne 527 bleiben 4. $P = \frac{4}{50} = 0{,}08 = 8\,\%$, also weniger als 10\,\%.}{}
\end{pfloesung}
\begin{pfloesung}{}
\lz{20b)}{2 rote Felder}{$\frac{1}{4} \cdot 8 = 2$ Felder}{}
\end{pfloesung}
\begin{pfloesung}{}
\lz{21.}{8 weiße Kugeln}{Die 4 schwarzen sind $\frac{1}{3}$ aller Kugeln, also alle $3 \cdot 4 = 12$ Kugeln; weiße: $12 - 4 = 8$}{}
\end{pfloesung}
\begin{pfloesung}{}
\lz{22.}{$\approx 0{,}173$ gegen $\approx 0{,}167$: passt gut}{relative Häufigkeit $\frac{52}{300} \approx 0{,}173$; $P(6) = \frac{1}{6} \approx 0{,}167$. Fast gleich: Bei vielen Würfen liegt die relative Häufigkeit nahe an der Wahrscheinlichkeit.}{}
\end{pfloesung}
\begin{pfloesung}{}
\lz{23.}{Nein: etwa 60 nötig, rund 10 zu wenig}{Nein. $P = \frac{2}{16} = \frac{1}{8}$; erwartet $\frac{1}{8} \cdot 480 = 60$ Teddys; es fehlen etwa $60 - 50 = 10$.}{}
\end{pfloesung}
\begin{pfloesung}{}
\lz{24.}{$\frac{1}{2}$}{4 von 8 Feldern sind rot, $\frac{4}{8} = \frac{1}{2}$. Nicht $\frac{1}{3}$ – es zählen die Felder, nicht die Farben.}{}
\end{pfloesung}
\begin{pfloesung}{}
\lz{25.}{$\frac{2}{15} \approx 13{,}3\,\%$}{alle $104 + 16 = 120$, günstig 16: $P = \frac{16}{120} = \frac{2}{15} \approx 13{,}3\,\%$}{}
\end{pfloesung}
\begin{pfloesung}{}
\lz{26.}{$P = \frac{3}{4}$}{über 9: 10, 11, 12 – das sind 3; $P = 1 - \frac{3}{12} = \frac{9}{12} = \frac{3}{4}$}{}
\end{pfloesung}
\begin{pfloesung}{}
\lz{27.}{Nein: $\frac{3}{12} = \frac{1}{4}$}{Nein. Es sind nur noch $15 - 3 = 12$ Bonbons, davon 3 mit Zitrone: $P = \frac{3}{12} = \frac{1}{4}$.}{}
\end{pfloesung}
\begin{pfloesung}{}
\lz{28.}{3 Felder; etwa 75 nötig – 70 reichen nicht}{$0{,}25 \cdot 12 = 3$ Preisfelder; erwartet $0{,}25 \cdot 300 = 75$ Preise. Nein, 70 reichen wahrscheinlich nicht; es fehlen etwa $75 - 70 = 5$.}{}
\end{pfloesung}
\end{document}
